\documentclass[12pt,a4paper]{article}

\usepackage{polyglossia}
\setmainlanguage{french}

\usepackage{amsmath, amssymb}

\usepackage{graphicx}
\usepackage{caption}

\usepackage{siunitx} % Pour les unités
\usepackage{myunit}


\usepackage{geometry}
\geometry{top=2cm, bottom=2.5cm, left=3cm, right=3cm}
\setlength{\parindent}{0pt}

\usepackage{enumitem}

\usepackage{hyperref}

\usepackage{xcolor}

\usepackage{multicol}
\setlength{\columnseprule}{1pt}
\def\columnseprulecolor{\color{gray}}

\usepackage{tcolorbox}


\usepackage{import}
\usepackage{xifthen}
\usepackage{pdfpages}
\usepackage{transparent}
\usepackage{float}

\newcommand{\incfig}[1]{%
	\def\svgwidth{\columnwidth}
	\import{./figures/}{#1.pdf_tex}
}


\newcommand{\RomanNum}[1] %pour les chiffre romain
{\MakeUppercase{\romannumeral #1}}



\usepackage{stanli}

\usepackage{pdflscape}

\usepackage{multirow}
\usepackage{makecell}





\title{TD 2 Bâtiment}
\author{}
\date{\today}

\begin{document}
	\counterwithin{section}{part} % pour que \section{} recommence de 1 aprés \part{}
	
	\maketitle
	
	\tableofcontents
	
	\newpage
	
	
	
	\part{Mur de soutènement}
	
	\section{Exo}
	
	
	Soit a pré dimensionner et vérifier la stabilité au renversement et au glissement du mur de soutènement suivant :
	
	
	\begin{center}
		
		\begin{tikzpicture}
			
			
			\draw (0,0) coordinate (a) -- (7,0) coordinate (b)
			-- (7,-1) coordinate (c) -- (8,-1) coordinate (d)
			-- (8,1) coordinate (e) -- (5,1) coordinate (f) 
			-- (5,6) coordinate (g) -- (4,6) coordinate (h)
			-- (3,1) coordinate (i) -- (0,1) coordinate (j) 
			-- (a);
			
			\point{G}{8}{6};
			
			\lineload{1}{g}{G};
			\notation{1}{G}{Q}[above left = 17mm];
			
			\dimensioning{2}{a}{h}{-.5}[$H = 8\meterl$];
			
			
		\end{tikzpicture}
		
	\end{center}
	
	
	
	\begin{itemize}
		\item Hauteur : $H = 8\meterl$
		\item Surcharge : $Q = 5\knm$
		\item $\varphi = \varPhi = \phi = 25\degree$
		\item $\gamma_{sol} = 1800\kgmc = 18\knmc$
		\item $\gamma_{ba} = 25\knmc$
	\end{itemize}
	
	
	La vérification au glissement se sera par la méthode du tiers central.
	
	
	\begin{tcolorbox}[colback=gray!20, colframe=black, title=Remarque]
		Le calcul se fera sur 1 mètre linéaire.
	\end{tcolorbox}
	
	
	
	
	
	
	
	\section{Pré Dimensionnement}
	
	\begin{center}
		
		\begin{tikzpicture}
			
			\draw (0,0) coordinate (a) -- (7,0) coordinate (b)
			-- (7,-1) coordinate (c) -- (8,-1) coordinate (d)
			-- (8,1) coordinate (e) -- (5,1) coordinate (f) 
			-- (5,6) coordinate (g) -- (4,6) coordinate (h)
			-- (3,1) coordinate (i) -- (0,1) coordinate (j) 
			-- (a);
			
			\dimensioning{1}{a}{d}{-2.5}[B];
			\dimensioning{1}{c}{d}{-1.8}[L];
			
			\dimensioning{1}{i}{f}{-1}[a'];
			\dimensioning{1}{h}{g}{6.5}[a];
			
			\dimensioning{1}{j}{i}{1.5}[$b_{1}$];
			\dimensioning{1}{f}{e}{1.5}[$b_{2}$];
			
			
			
			\dimensioning{2}{e}{b}{8.5}[h];
			\dimensioning{2}{b}{d}{8.5}[h'];
			
			\dimensioning{2}{h}{a}{-1}[H];
			
			
			
		\end{tikzpicture}
		
	\end{center}
	
	
	
	\subsection{Le voile}
	
	\[ a \geq \dfrac{H}{24} = \dfrac{8}{24} = 0,33\meterl \]
	
	\[ a' \geq \dfrac{H}{12} = 0,66\meterl \]
	
	On prend : \fbox{$a = 35\cml$ et $a' = 70\cml$}
	
	
	
	\subsection{La semelle}
	
	\subsubsection{La hauteur}
	
	\[ h \geq \dfrac{H}{12} = 0,66\meterl \]
	
	\[ \fbox{$h = 70\cml$} \]
	
	\subsubsection{La largeur}
	
	\[ \dfrac{H}{2} < B \leq \dfrac{2H}{3} \]
	
	\[ 4\meterl < B \leq 5,33\meterl \]
	
	\[ \fbox{$B = 4,7\meterl$} \]
	
	\subsection{Le patin}
	
	\[ b_{1} \geq \dfrac{B}{3} = \dfrac{4,7}{3} = 1,56\meterl \]
	
	\[ \fbox{$b_{1} = 1,6\meterl$} \]
	
	\subsection{Le talon}
	
	\[ b_{2} = B - (b_{1} + a') = 4,7 - (1,6 + 0,7) \]
	
	\[ \fbox{$b_{2} = 2,4\meterl$} \]
	
	\subsubsection{La bêche}
	
	\[ \fbox{$h' = 0,7\meterl$} \]
	\[ \fbox{$L = 0,7\meterl$} \]
	
	
	
		\begin{center}
		
		\begin{tikzpicture}[scale=.8]
			
			\draw (0,0) coordinate (a) -- (7,0) coordinate (b)
			-- (7,-1) coordinate (c) -- (8,-1) coordinate (d)
			-- (8,1) coordinate (e) -- (5,1) coordinate (f) 
			-- (5,6) coordinate (g) -- (4,6) coordinate (h)
			-- (3,1) coordinate (i) -- (0,1) coordinate (j) 
			-- (a);
			
			\dimensioning{1}{a}{d}{-2.5}[4,7];
			\dimensioning{1}{c}{d}{-1.8}[0,7];
			
			\dimensioning{1}{i}{f}{-1}[0,7];
			\dimensioning{1}{h}{g}{6.5}[0,35];
			
			\dimensioning{1}{j}{i}{1.5}[1,6];
			\dimensioning{1}{f}{e}{1.5}[2,4];
			
			
			
			\dimensioning{2}{e}{b}{8.5}[0,7];
			\dimensioning{2}{b}{d}{8.5}[0,7];
			
			\dimensioning{2}{h}{a}{-1}[8];
			
			\node at (a) [circle,draw=red!50,fill=red!80] {C};
			
		\end{tikzpicture}
		
	\end{center}
	
	
	\section{Évaluations des charges}
	
	\newpage
		
		
	\begin{landscape}
	
		\begin{table}[h!]
			\begin{center}
				\begin{tabular}{l|l|l|l|l}
					
					Types de charges & partie & Force & Distance vers c & Moment/c ($G \cdot X$) \\ \hline
					
					
					
					
					\multirow{5}{*}{Charges verticales} &
					
						\begin{tikzpicture}[scale=.3]
							
							\draw (0,0) -- (3,0) -- (3,6) -- (2,6) -- (0,0);
							
							\draw[<-,very thick] (2,2) coordinate(a) -- (2,4) node[above]{$G_{voile}$};
							
							\point{b}{-1.5}{0};
							
							\node at (b) [circle,draw=red!50,fill=red!80,] {C};
							\dimensioning{1}{a}{b}{-.5}[\small X];
							
						
						\end{tikzpicture} & $G_{voile} = 25 \cdot 0,35 \cdot 0,7 \cdot 1 = 95,81\kn $ & $X_{voile} = 2,03\meterl $ & $M_{voile} = 194,5\knxm$ \\ \cline{2-5}
						
						
						
						
						
						& \begin{tikzpicture}[scale=.3]
							
							\draw (0,0) coordinate(a) node[circle,draw=red!50,fill=red!80,] {C} -- (7,0) -- (7,1) -- (0,1) -- (0,0);
							
							\draw[<-,very thick] (3.5,.5) coordinate(b) -- (3.5,2) node[above]{$G_{semelle}$};
							
							\dimensioning{1}{a}{b}{-.5}[X];
							
							
						\end{tikzpicture} & $G_{semelle} = 25 \cdot 4,7 \cdot 0,7 \cdot 1 = 82,25\kn$ & $X_{semelle} = \dfrac{4,7}{2} = 2,35\meterl$ & $M_{semelle} = 193,28\knxm$ \\ \cline{2-5}
						
						
						
						
						
						
						& \begin{tikzpicture}[scale=.5]
							
							\draw(0,0) rectangle (1,1);
							
							\draw[<-,very thick] (.5,.5) -- (.5,1.5) node[above]{$G_{beche}$} coordinate (a);
							
							\point{b}{-2}{0};
							\node at (b) [circle,draw=red!50,fill=red!80,below] {C};
							\dimensioning{1}{a}{b}{-.5}[X];
							
							
						\end{tikzpicture}& $G_{beche} = 25 \cdot 0,7 \cdot 0,7 \cdot 1 = 12,25\kn$ & $X_{beche} = 4,7 - \dfrac{0,7}{2} = 4,35\meterl$ & $M_{beche} = 53,28\knxm$ \\ \cline{2-5}
						
						
						
						
						
						
						& \begin{tikzpicture}[scale=.3]
							\draw (0,0) coordinate (a) -- (7,0) coordinate (b)
							-- (7,-1) coordinate (c) -- (8,-1) coordinate (d)
							-- (8,1) coordinate (e) -- (5,1) coordinate (f) 
							-- (5,6) coordinate (g) -- (4,6) coordinate (h)
							-- (3,1) coordinate (i) -- (0,1) coordinate (j) 
							-- (a);
							
							\draw[brown,very thick] (f) rectangle (8,6);
							\draw[<-,very thick] (6.5,3) -- (6.5,6) coordinate(F) node[above]{$G_{sol}$};
							
							\dimensioning{1}{a}{F}{-1}[X];
							
						\end{tikzpicture} & $G_{sol} = 18 \cdot 2,4 \cdot 7,3 \cdot 1 = 315,36\kn$ & $X_{sol} = 4,7 - \dfrac{2,4}{2} = 3,5\meterl$ & $M_{sol} = 1103,76\knxm$ \\ \cline{2-5}
						
						
						
						
						
						
						& \begin{tikzpicture}
							
							\point{a}{0}{0};
							\point{b}{2}{0};
							
							\point{c}{-1}{0};
							
							\lineload{1}{a}{b};
							
							\draw[<-,very thick] (1,0) coordinate(F) -- (1,2) node[above]{$G_{surcharge}$};
							
							\node at (c) [circle,draw=red!50,fill=red!80,below] {C};
							
							\dimensioning{1}{F}{c}{-.5}[X];
						\end{tikzpicture} & $G_{surcharge} = 5 \cdot 2,4 = 12\kn$ & $X_{surcharge} = X_{sol} = 3,5\meterl$ & $M_{surcharge} = 42\knxm$ \\ \hline
						
						
						
						
						
						Charges horizontales & \makecell{Poussée du sol \\ et de la surcharge} \\
						& \begin{tikzpicture}[scale=.3]
								\draw (0,0) coordinate (a) -- (5,0) coordinate (F) -- (7,0) coordinate (b)
							-- (7,-1) coordinate (c) -- (8,-1) coordinate (d)
							-- (8,1) coordinate (e) -- (5,1) coordinate (f) 
							-- (5,6) coordinate (g) -- (4,6) coordinate (h)
							-- (3,1) coordinate (i) -- (0,1) coordinate (j) 
							-- (a);
							
							\lineload{1}{g}{F}[2][4];
							
							\draw[<-,very thick,red] (5,2) -- (11,2) coordinate (A) node[above,right]{$F_{p}$};
							
							\dimensioning{2}{A}{a}{11.5}[Y];
							
						\end{tikzpicture} & \makecell{$F_{p} = 250,52\kn$ \\ voir la page suivante pour les détails} & $Y = 2,75\meterl$ & $M_{r} = 688,93\knxm$
						
						
						
					
					
					
					
				\end{tabular}
			\end{center}
		\end{table}
		
		
	\end{landscape}		
	
	
	
	\newpage
	
	\subsection{Charges horizontales}
	
	Nous avons la poussée du sol + surcharge : 
	
	\[ \sigma_{h} = K \cdot \sigma_{v} = K \cdot (\gamma_{sol} \cdot x \cdot 1ml + Q) \]
	
	\[ K = \tan^{2}(\dfrac{\pi}{4} - \dfrac{\varphi}{2}) = \tan^{2}(\dfrac{180}{4} - \dfrac{25}{2}) = 0,406 \]
	
	
	\[ \sigma_{h} = 0,406 \cdot (18 \cdot x + 5) \rightarrow \begin{cases}
		x = 0 \rightarrow \sigma_{h} = 0,406 \cdot 5 = 2,03\knms \\ \\
		x = 8\meterl \rightarrow \sigma_{h} = 0,406 \cdot (18 \cdot 8 + 5) = 60,5\knms
	\end{cases} \]
	
	
	\begin{center}
		\begin{tikzpicture}
			\draw (0,0) coordinate (a) -- (5,0) coordinate (F) node[below]{$60,5\knms$} -- (7,0) coordinate (b)
			-- (7,-1) coordinate (c) -- (8,-1) coordinate (d)
			-- (8,1) coordinate (e) -- (5,1) coordinate (f) 
			-- (5,6) coordinate (g) node[above]{2,03\knms} -- (4,6) coordinate (h)
			-- (3,1) coordinate (i) -- (0,1) coordinate (j) 
			-- (a);
			
			\lineload{1}{g}{F}[2][4];
			
			\draw[<-,very thick,red] (5,2) -- (11,2) coordinate (A) node[above,right]{$F_{p}$};
			
			\dimensioning{2}{A}{a}{12}[Y];
			\dimensioning{2}{a}{h}{-.2}[8m];
			
		\end{tikzpicture}
	\end{center}
	
	
	\[ F_{p} = \text{Surface du trapéze} = \dfrac{(2,03 + 60,5) \cdot 8}{2} = 250,12\kn \]
	
	\[ Y = \dfrac{2,03 \cdot 8 \cdot 4 + \dfrac{8 \cdot 58,47}{2} \cdot \dfrac{8}{3}}{250,12} = 2,75\meterl \]
	
	\begin{tcolorbox}[colback=gray!20, colframe=black, title=Remarque]
		On peut calculer la poussée du sol seul et celle de la surcharge seul.
	\end{tcolorbox}
	
	
	\section{Vérification de la stabilité}
	
	\[ \dfrac{\sum M_{s}}{\sum M_{r}} \geq 1,5 \]
	
	\[ \sum M_{s} = 194,5 + 193,28 + 53,28 + 1103,76 + 42  =  1586,82\knxm \]
	
	\[ \sum M_{r} = 688,93\knxm \]
	
	\[ \dfrac{1586,82}{688,93} = 2,3 > 1,5 \] 
	
	{\centering Condition vérifier le mur est stable vis a vis au renversement \par}
	
	\section{Vérification au glissement avec Tiers central}
	
	\paragraph{Force verticale} $\sum G = 517,67\kn = R_{v}$
	
	\begin{center}
		\begin{tikzpicture}
			\draw (0,0) coordinate (a) -- (7,0) coordinate (b)
			-- (7,-1) coordinate (c) -- (8,-1) coordinate (d)
			-- (8,1) coordinate (e) -- (5,1) coordinate (f) 
			-- (5,6) coordinate (g) -- (4,6) coordinate (h)
			-- (3,1) coordinate (i) -- (0,1) coordinate (j) 
			-- (a);
			
			\draw[<-,very thick,red] (6.5,3) -- (6.5,6) coordinate(F) node[above]{$R_{v}$};
			
			\dimensioning{1}{a}{F}{-1}[X];
			
		\end{tikzpicture}
	\end{center}
	
	La condition est la suivante : 
	\[ X \in [\dfrac{B}{2}\,;\,\dfrac{2B}{3}] \]
	
	Nous avons : moments stabilisateurs : $\sum M_{s} = R_{v} \cdot x$
	
	Donc : \[ X = \dfrac{\sum M_{s}}{R_{v}} = \dfrac{1586,82}{517,67} = 3,06\meterl \]
	
	\[ X = 3,06\meterl \in [2,35 \quad 3,13] \]
	
	{\centering Condition vérifier, le mur est stable vis a vis au glissement \par}
	
	
\end{document}